% 00_overview.tex - what the workshop covers, what you need, how to read it
\begin{goal}
\textbf{What you do.} You calibrate the groundwater model of Tutorial Creek against the streamflow at the outlet and the
heads at three observation wells, with the Dynamically Dimensioned Search (DDS) of the program OSTRICH. The observations
were made with known (``true'') parameters, so you can check what the calibration recovers. First you look at which
parameters the observations can see at all, and at a first attempt that went wrong. Then you calibrate the two parameters
that the observations see clearly and separately: the hydraulic conductivities of the valley sand and of the slope till.
Finally you validate the result on two years that the calibration has not seen. The set-up follows the official
Raven--OSTRICH tutorial (Demo~C1: \cmd{ostIn.txt}, a template, \cmd{Ost-RAVEN.sh} and \cmd{save\_best.sh}).

\textbf{What you need.} A Linux computer (x86\_64) with a terminal, or Windows with WSL. OSTRICH, which is not included
(Step~\ref{st:ostrich} says where to get it). A text editor, for example \cmd{nano}. Workshop~2 is helpful but not needed.

\textbf{Time.} About 30 minutes of work, plus about \V{t6:ddsmin}~minutes of computing for the calibration
(\V{t6:maxiter} model runs).
\end{goal}
